\section{DPM-Solver de orden superior}

\subsection{Aproximación de segundo orden: DPM-Solver-2}

La aproximación de primer orden asume que $\boldsymbol{\epsilon}_\theta$ es constante en todo el intervalo. La mejora natural consiste en utilizar una \textbf{expansión de Taylor de primer orden} (aproximación lineal), empleando las predicciones de ruido en dos puntos para ajustar la tendencia de cambio de $\boldsymbol{\epsilon}_\theta$.

Los pasos del algoritmo DPM-Solver-2 son los siguientes:

\textbf{Paso 1: Calcular el punto intermedio}

Dado el paso temporal actual $s$ y el paso temporal objetivo $t$, seleccionar un punto temporal intermedio $r$ (generalmente el punto medio del intervalo) y calcular el valor intermedio con DPM-Solver de primer orden:
$$\mathbf{u} = \frac{\alpha_r}{\alpha_s} \mathbf{x}_s - \sigma_r (e^{h_1} - 1) \boldsymbol{\epsilon}_\theta(\mathbf{x}_s, s)$$

Donde $h_1 = \lambda_r - \lambda_s$.

\textbf{Paso 2: Actualización utilizando la predicción de ruido en el punto intermedio}

Reevaluar la predicción de ruido $\boldsymbol{\epsilon}_\theta(\mathbf{u}, r)$ en el punto intermedio $r$ y luego realizar una aproximación por interpolación lineal usando la información de dos puntos para integrar:
$$\mathbf{x}_t = \frac{\alpha_t}{\alpha_s} \mathbf{x}_s - \sigma_t (e^{h} - 1) \boldsymbol{\epsilon}_\theta(\mathbf{x}_s, s) - \frac{\sigma_t}{2h_1}(e^{h} - 1 - h)\left[\boldsymbol{\epsilon}_\theta(\mathbf{u}, r) - \boldsymbol{\epsilon}_\theta(\mathbf{x}_s, s)\right]$$

Donde $h = \lambda_t - \lambda_s$.

\begin{importantbox}{La idea clave de la aproximación de segundo orden}
DPM-Solver-2 requiere 2 evaluaciones de red por paso (en $s$ y en $r$), pero logra precisión de segundo orden. A diferencia del método de Heun general, DPM-Solver-2 aprovecha la estructura semilineal de PF-ODE, lo que resulta en mayor precisión para el mismo número de NFE.
\end{importantbox}

\subsection{Aproximaciones de tercer orden y superiores}

Siguiendo el mismo razonamiento, se puede continuar con una expansión de Taylor de segundo orden para obtener DPM-Solver-3: seleccionar dos puntos intermedios en el intervalo y utilizar tres predicciones de ruido para construir una aproximación polinómica cuadrática del cambio de $\boldsymbol{\epsilon}_\theta$ en el espacio $\lambda$.

\begin{knowledgebox}{Estructura recursiva}
El aumento del orden en DPM-Solver tiene una estructura recursiva:
\begin{itemize}
    \item \textbf{DPM-Solver de orden $k$} requiere $k$ evaluaciones de red
    \item Primero calcular $k-1$ puntos intermedios con un DPM-Solver de orden inferior
    \item Utilizar todas las $k$ predicciones de ruido para construir una aproximación polinómica de grado $(k-1)$
    \item Sustituir el polinomio en la fórmula de integral y realizar la integración analítica
\end{itemize}
En la práctica, los órdenes 2 y 3 son suficientes. Más allá del tercer orden, las ganancias disminuyen porque el error de aproximación inherente a la red de predicción de ruido se convierte en un cuello de botella.
\end{knowledgebox}

\subsection{Comparación de DPM-Solver por órdenes}

\begin{table}[H]
\centering
\begin{tabular}{lccc}
\toprule
\textbf{Método} & \textbf{NFE por paso} & \textbf{Orden de precisión} & \textbf{Observaciones} \\
\midrule
DDPM & 1 & Aprox. 0.5 orden & Euler-Maruyama discreción simple \\
DDIM / DPM-Solver-1 & 1 & 1 orden & Integral exponencial + aproximación de ruido de orden cero \\
DPM-Solver-2 & 2 & 2 orden & Integral exponencial + aproximación de ruido de primer orden \\
DPM-Solver-3 & 3 & 3 orden & Integral exponencial + aproximación de ruido de segundo orden \\
\bottomrule
\end{tabular}
\caption{Comparación de las características de los métodos DPM-Solver por órdenes}
\end{table}

\subsection{Resumen de la sección}

Los DPM-Solver de orden superior logran mayor precisión al realizar aproximaciones polinómicas de alto orden de las predicciones de ruido en el espacio $\lambda$, aumentando solo ligeramente el costo de evaluaciones de red por paso. Esta construcción recursiva hace que el avance del primer al tercer orden sea conceptualmente muy natural.
